% ============================================================================
%                           LaTeX 最简文档模板
%                        使用 ctexart 文档类（中文支持）
% ============================================================================

% ----------------------------------------------------------------------------
% 第一部分：文档类声明
% ----------------------------------------------------------------------------
% \documentclass[选项]{文档类}
% - [UTF8]：指定文档编码为 UTF-8，支持中文
% - {ctexart}：中文文章文档类，基于 article 扩展
\documentclass[UTF8]{ctexart}

% ----------------------------------------------------------------------------
% 第二部分：宏包加载区（导言区）
% ----------------------------------------------------------------------------
% 所有 \usepackage 命令必须放在 \documentclass 之后、\begin{document} 之前

% amsmath：提供增强的数学公式支持
\usepackage{amsmath}

% xcolor：提供颜色支持
\usepackage{xcolor}

% geometry：页面布局设置
\usepackage{geometry}
\geometry{a4paper, scale=0.8}  % A4纸张，页面比例0.8

% ----------------------------------------------------------------------------
% 第三部分：文档元信息设置
% ----------------------------------------------------------------------------
% \title{标题}：设置文档标题
\title{文档标题}

% \author{作者}：设置作者姓名
\author{作者姓名}

% \date{日期}：设置日期
% - \today：自动生成当前日期
% - 也可手动指定，如 \date{2024年1月1日}
\date{\today}

% ----------------------------------------------------------------------------
% 第四部分：正文区
% ----------------------------------------------------------------------------
% \begin{document} ... \end{document}：文档正文开始和结束标记
% 所有可见内容必须放在这两个命令之间
\begin{document}

% 生成标题
% \maketitle：根据上方设置的标题、作者、日期生成标题区域
\maketitle

% 生成目录
% \tableofcontents：自动生成目录（可选）
% \tableofcontents

% ----------------------------------------------------------------------------
% 正文内容
% ----------------------------------------------------------------------------

% 章节命令（ctexart 提供中文友好的章节）
% \section{一级标题}
\section{引言}

% 段落：LaTeX 中空行表示段落分隔
这是一个简单的段落。LaTeX 会自动处理段落缩进和换行。

空一行开始新段落。你可以在这里输入更多内容。

% \subsection{二级标题}
\subsection{子标题}

这是子标题下的段落内容。

% ----------------------------------------------------------------------------
% 常用元素示例
% ----------------------------------------------------------------------------

\section{常用元素}

\subsection{列表}

% 无序列表
\begin{itemize}
    \item 第一项
    \item 第二项
    \item 第三项
\end{itemize}

% 有序列表
\begin{enumerate}
    \item 第一项
    \item 第二项
    \item 第三项
\end{enumerate}

\subsection{数学公式}

% 行内公式：使用 $...$ 包围
爱因斯坦质能方程：$E=mc^2$

% 行间公式：使用 \[...\] 或 equation 环境
\[
    a^2 + b^2 = c^2
\]

\subsection{强调}

% 文本强调方式
\textbf{粗体文本}

\textit{斜体文本}

\underline{下划线文本}

% ----------------------------------------------------------------------------
% 文档结束
% ----------------------------------------------------------------------------
\end{document}

% ============================================================================
% 模板说明：
% 
% 编译方法：
%   xelatex 文档名.tex  或  pdflatex 文档名.tex
% 
% 文件结构：
%   1. 文档类声明（必须放在第一行）
%   2. 宏包加载（在导言区）
%   3. 元信息设置（标题、作者、日期）
%   4. 正文区（\begin{document} 到 \end{document}）
% 
% 扩展建议：
%   - 添加 \usepackage{graphicx} 支持图片
%   - 添加 \usepackage{booktabs} 支持精美表格
%   - 添加 \usepackage{hyperref} 支持超链接
% ============================================================================
